\begin{tikzpicture}[
    >=Latex,
    scale=0.95,
    every node/.style={font=\small},
]
    % ===========================================
    % PANEL A: Ohne Mittenversatz (links)
    % ===========================================
    \begin{scope}[xshift=0cm]
        % Titel
        \node[font=\bfseries, anchor=south] at (1, 5)
            {(a) An Schichtgrenzen};

        % Aufbauachse a
        \draw[->, thick] (-0.7, -0.3) -- (-0.7, 4.3)
            node[above, font=\scriptsize] {$\vec{a}$};

        % Bauteil (Querschnitt) - Hoehe 3.7 t
        \fill[gray!15] (0, 0) rectangle (1.5, 3.7);
        \draw[thick] (0, 0) rectangle (1.5, 3.7);

        % a_min / a_max Markierungen
        \draw[dotted, thin] (0, 0) -- (-0.85, 0);
        \node[anchor=east, font=\scriptsize] at (-0.85, 0) {$a_{\min}$};
        \draw[dotted, thin] (0, 3.7) -- (-0.85, 3.7);
        \node[anchor=east, font=\scriptsize] at (-0.85, 3.7) {$a_{\max}$};

        % Schichtgrenzen-Hilfslinien
        \foreach \i in {1,2,3} {
            \draw[gray!40, dotted, very thin] (-0.3, \i) -- (2.3, \i);
        }

        % Schnittebenen (an Schichtgrenzen)
        \foreach \y/\i in {0/0, 1/1, 2/2, 3/3} {
            \draw[blue!70!black, very thick] (-0.1, \y) -- (1.6, \y);
            \node[blue!70!black, anchor=west, font=\scriptsize]
                at (1.65, \y) {$a_{\i}$};
        }

        % Nicht erfasster Bereich
        \fill[red!25] (0, 3) rectangle (1.5, 3.7);
        \draw[red!70!black, thick] (0, 3) rectangle (1.5, 3.7);

        % Annotation des fehlenden Fragments
        \node[red!80!black, anchor=south, font=\scriptsize, align=center]
            at (0.75, 4.0) {fehlende Schicht\\($0{,}7\,t$)};

        % Formel
        \node[font=\scriptsize, anchor=north] at (0.75, -0.6)
            {$a_i = a_{\min} + i \cdot t$};
    \end{scope}

    % ===========================================
    % PANEL B: Mit Mittenversatz (rechts)
    % ===========================================
    \begin{scope}[xshift=6cm]
        % Titel
        \node[font=\bfseries, anchor=south] at (1, 5)
            {(b) Mit Mittenversatz};

        % Aufbauachse a
        \draw[->, thick] (-0.7, -0.3) -- (-0.7, 4.3)
            node[above, font=\scriptsize] {$\vec{a}$};

        % Bauteil
        \fill[gray!15] (0, 0) rectangle (1.5, 3.7);
        \draw[thick] (0, 0) rectangle (1.5, 3.7);

        % a_min / a_max
        \draw[dotted, thin] (0, 0) -- (-0.85, 0);
        \node[anchor=east, font=\scriptsize] at (-0.85, 0) {$a_{\min}$};
        \draw[dotted, thin] (0, 3.7) -- (-0.85, 3.7);
        \node[anchor=east, font=\scriptsize] at (-0.85, 3.7) {$a_{\max}$};

        % Schichtgrenzen-Hilfslinien
        \foreach \i in {0,1,2,3,4} {
            \draw[gray!40, dotted, very thin] (-0.3, \i) -- (2.3, \i);
        }

        % Schnittebenen (mittig)
        \foreach \y/\i in {0.5/0, 1.5/1, 2.5/2, 3.5/3} {
            \draw[green!50!black, very thick] (-0.1, \y) -- (1.6, \y);
            \node[green!50!black, anchor=west, font=\scriptsize]
                at (1.65, \y) {$a_{\i}$};
        }

        % Bewusst akzeptierter Restbereich (klein)
        \fill[gray!40] (0, 3.5) rectangle (1.5, 3.7);

        % Annotation
        \node[gray!60!black, anchor=south, font=\scriptsize, align=center]
            at (0.75, 4.0) {$<\tfrac{t}{2}$ (akzeptiert)};

        % t/2 und t Bemassung innerhalb des Bauteils
        \draw[<->, thin, blue!70!black] (0.15, 0) -- (0.15, 0.5);
        \node[anchor=west, font=\scriptsize, blue!70!black]
            at (0.18, 0.25) {$\tfrac{t}{2}$};
        \draw[<->, thin, blue!70!black] (0.15, 0.5) -- (0.15, 1.5);
        \node[anchor=west, font=\scriptsize, blue!70!black]
            at (0.18, 1.0) {$t$};

        % Formel
        \node[font=\scriptsize, anchor=north] at (0.75, -0.6)
            {$a_i = a_{\min} + (i+\tfrac{1}{2}) \cdot t$};
    \end{scope}
\end{tikzpicture}
